\documentclass[10.5pt,a4paper]{ctexart}

% 使用 XeLaTeX 编译。
\usepackage{geometry}
\usepackage{graphicx}
\usepackage{float}
\usepackage{booktabs}
\usepackage{array}
\usepackage{makecell}
\usepackage{enumitem}
\usepackage{amsmath}
\usepackage{amssymb}
\usepackage{caption}
\usepackage{subcaption}
\usepackage{hyperref}

\geometry{top=2.54cm,bottom=2.54cm,left=2.54cm,right=2.0cm}
\setlength{\parindent}{2em}
\setlength{\parskip}{0.25em}
\linespread{1.2}
\setlist[enumerate]{label=(\arabic*),leftmargin=2.5em}
\setlist[itemize]{leftmargin=2.5em}
\hypersetup{
  colorlinks=true,
  linkcolor=blue,
  urlcolor=blue,
  citecolor=blue
}

\newcolumntype{C}[1]{>{\centering\arraybackslash}p{#1}}
\newcolumntype{L}[1]{>{\raggedright\arraybackslash}p{#1}}
\newcommand{\blank}[1]{\underline{\makebox[#1]{}}}

% ===== 报告信息 =====
\newcommand{\reportCourse}{电工电子学实验}
\newcommand{\reportTitle}{实验报告}
\newcommand{\experimentName}{}
\newcommand{\experimentType}{}
\newcommand{\experimentDate}{}
\newcommand{\experimentPlace}{}
\newcommand{\studentMajor}{}
\newcommand{\studentName}{}
\newcommand{\studentId}{}
\newcommand{\partnerName}{}
\newcommand{\instructorName}{}
\newcommand{\score}{}

\title{\reportCourse\\\reportTitle}
\author{}
\date{}

\begin{document}

\begin{center}
  {\zihao{2}\bfseries \reportCourse\quad \reportTitle}
\end{center}

\begin{table}[H]
  \centering
  \renewcommand{\arraystretch}{1.45}
  \begin{tabular}{L{2.2cm}L{5.0cm}L{2.2cm}L{5.0cm}}
    实验名称 & \experimentName & 实验类型 & \experimentType \\
    实验日期 & \experimentDate & 实验地点 & \experimentPlace \\
    专业 & \studentMajor & 姓名 & \studentName \\
    学号 & \studentId & 同组学生 & \partnerName \\
    指导老师 & \instructorName & 成绩 & \score \\
  \end{tabular}
\end{table}

\section{实验目的和要求}
\begin{enumerate}
  \item 写明本次实验需要验证、测量或掌握的核心内容。
  \item 将目的按知识点或实验任务拆分，避免只写笼统表述。
\end{enumerate}

\section{实验仪器与设备}
\begin{table}[H]
  \centering
  \caption{主要仪器与元器件}
  \begin{tabular}{C{1.2cm}L{6.0cm}L{5.0cm}}
    \toprule
    序号 & 名称 & 型号或参数 \\
    \midrule
    1 &  &  \\
    2 &  &  \\
    3 &  &  \\
    \bottomrule
  \end{tabular}
\end{table}

\section{实验原理}
概述本实验涉及的基本定律、电路模型、关键公式和理论预期。

\subsection{基本原理}
例如：
\begin{equation}
  y = kx + b
\end{equation}
其中，各符号的含义、单位和适用条件应在正文中说明。

\subsection{理论分析}
可结合电路图、波形图或等效模型说明测量量之间的关系，并指出后续数据处理所需的计算式。

\section{实验内容与步骤}
\begin{enumerate}
  \item 按实验指导书连接电路或搭建仿真模型。
  \item 检查仪器量程、接线极性和公共端连接。
  \item 改变实验条件，记录测量数据或波形。
  \item 关闭电源，整理数据并进行误差分析。
\end{enumerate}

\section{实验数据记录与处理}
\subsection{原始数据}
\begin{table}[H]
  \centering
  \caption{实验数据记录表}
  \begin{tabular}{C{2cm}C{2.5cm}C{2.5cm}C{2.5cm}C{2.5cm}}
    \toprule
    序号 & 测量条件 & 测量量 1 & 测量量 2 & 备注 \\
    \midrule
    1 &  &  &  &  \\
    2 &  &  &  &  \\
    3 &  &  &  &  \\
    \bottomrule
  \end{tabular}
\end{table}

\subsection{数据处理}
写出计算过程、单位换算、误差或相对偏差：
\begin{equation}
  \delta = \frac{|x_{\mathrm{meas}} - x_{\mathrm{theory}}|}{x_{\mathrm{theory}}}\times 100\%
\end{equation}

\section{实验波形与图像}
插入实验电路、仿真结果或示波器波形。图片文件建议放在同级目录的 \texttt{figures/} 文件夹中。

\begin{figure}[H]
  \centering
  % \includegraphics[width=0.65\textwidth]{figures/example.png}
  \fbox{\parbox[c][4cm][c]{0.65\textwidth}{\centering 图片占位}}
  \caption{实验图像或波形}
  \label{fig:example}
\end{figure}

\section{实验结果与分析}
\begin{enumerate}
  \item 对照理论值和实测值，说明实验现象是否符合预期。
  \item 分析主要误差来源，例如仪表内阻、读数误差、元件容差、接触电阻、波形失真等。
  \item 对异常数据或波形给出原因判断，而不是只重复表格中的数值。
\end{enumerate}

\section{思考题}
\begin{enumerate}
  \item 题目一：写出题干或要点。
  \item 解答：结合实验原理和数据作答。
\end{enumerate}

\section{讨论与心得}
总结本次实验的关键收获、操作注意事项和后续改进方向。若涉及工程应用，可说明该实验结论在实际电路或设备中的意义。

\end{document}
